\chapter{总结与展望}

\section{工作总结}
本文针对机器人灵巧遥操作中运动重定向的平滑性、精准性以及跨平台通用性问题，提出了基于时空 Transformer 的运动重定向框架 TransRetarget。通过理论研究、算法设计与多平台实验验证，本文的主要工作总结如下：

构建了高效的时空特征提取架构：借鉴 PoseFormer 的分级建模思想，设计了专门针对手部三维骨骼序列的空间编码器与时间编码器。空间编码器利用自注意力机制捕获单帧内指尖与掌心间的复杂拓扑关系；时间编码器则通过对连续多帧数据的时域建模，提取运动趋势特征，有效解决了传统方法中常见的动作抖动与不连贯问题。

设计了多维度物理约束损失函数：为了平衡运动的保真度与物理安全性，本文设计了包含手指摆动、指尖位置、指尖间距以及碰撞惩罚在内的复合损失函数。该设计确保了重定向结果不仅在空间坐标上精准还原人手意图，且在运动学配置上符合机器人硬件的物理限位，显著提升了多指协作抓取的稳定性。

实现了跨平台的灵巧手适配能力：通过引入基于 URDF 解析的通用运动学模块，本文算法成功在 6 自由度的 Inspire 手、17 自由度的 Linker 手以及 23 自由度的 Shadow 手上完成了部署。实验结果表明，该框架能够自动适配不同自由度分布的机器人硬件，具有极强的泛化性能。

验证了现实场景的实用价值：通过单机现实世界实验，验证了算法在处理 6 种不同几何特性物体抓取任务时的鲁棒性。在 60 FPS 的实时处理速度下，系统能够平滑、精准地复现人类操作员的精细捏取与握持动作，为低延迟遥操作任务提供了可靠的技术支撑。

\section{研究展望}
尽管本文在运动重定向的平滑度与跨平台迁移方面取得了阶段性成果，但面对更加复杂、多变的真实交互场景，仍有以下方向值得进一步探索：

引入频域优化提升长序列处理能力：目前的时域自注意力机制在处理极长序列时计算开销较大。未来可尝试引入频域特征提取（如离散余弦变换等）来捕捉周期性或快速变化的运动特征，在保持鲁棒性的同时进一步降低计算复杂度。

增强力反馈与物理交互感知：目前的重定向主要基于位置映射。后续研究可考虑将末端执行器的力矩信息或触觉反馈融入损失函数，实现具备“触感”的运动重定向，从而在处理易碎物体或需要恒力交互的任务中表现得更加出色。

大模型先验知识的注入：探索将预训练的视觉-语言大模型或通用的具身智能大模型作为特征提取的先验，使重定向系统能够理解更高层次的动作语义（如“递给”、“旋转”），从而在传感器数据缺失或遮挡严重的情况下实现智能化的动作补全。

手臂和手协同的重定向研究：目前的重点在于灵巧手部。未来将尝试将该框架扩展至全臂与灵巧手协同的重定向任务中，解决人体手臂与机器人机械臂之间的工作空间映射问题，实现更加自然、直观的全身遥操作体验。

